\documentclass[12pt]{scrartcl}

\usepackage[    % set up dimensions
	a6paper, 
	landscape,
	left=1.5cm,   
	right=1.5cm,  
	top=1.5cm,
	right=1.5cm, 
	showframe,
	nomarginpar,  % this removes the right margin, centering the text in the middle
	noheadfoot,  % this removes the head and foot margin
	]{geometry}

%\pagestyle{empty}

\usepackage{graphicx}
\usepackage{tikz}

% Fonts and encodings
\usepackage[nf]{coelacanth}
\usepackage[T1]{fontenc}

\usepackage{longtable} % Um Tabellen zu machen, die mehrere Seiten lang sind und einen passenden Seitenumbruch haben.
%\usepackage[singlespacing]{setspace} % Zeilenabstand
\usepackage{nicefrac} % für besser aussehende Brüche

%\usepackage[object=vectorian]{pgfornament}

\begin{document}
	\pagestyle{empty}
	
	\begin{tikzpicture}[remember picture, overlay]
		% pgfornament does give very simple and nice vectorian graphics for tikz.
		% The problem is that the lines are too thick and figure 41 is also slightly cropped.
		% pst-vectorian works very well.
		% we're now using the original files from vectorian exported as eps file
		
		% DEBUGGING grid
		%\draw[line width=0.75pt, color=teal, step=.5cm, xshift=0.08cm, yshift=0.43cm] (current page.south west) grid ++(14.85cm,10.5cm);
		% DEBUGGING corners
		%\node [circle, fill=red, minimum size=1pt, xshift=0.0cm, yshift=-0.0cm, scale=0.2] at (current page.north west) {};
		%\node [circle, fill=red, minimum size=1pt, xshift=0.5cm, yshift=-0.5cm, scale=0.2] at (current page.north west) {};
		
		
		\node[anchor=north west, xshift=0.5cm, yshift=-0.5cm, inner sep=0] at (current page.north west){\includegraphics[scale=0.8, angle=0, origin=c]{img/corner}};
		\node[anchor=north east, xshift=-0.5cm, yshift=-0.5cm, inner sep=0] at (current page.north east){\includegraphics[scale=0.8, angle=0, origin=c]{img/corner2}};
		\node[anchor=south west, xshift=0.5cm, yshift=0.5cm, inner sep=0] at (current page.south west){\includegraphics[scale=0.8, angle=180, origin=c]{img/corner2}};
		\node[anchor=south east, xshift=-0.5cm, yshift=0.5cm, inner sep=0] at (current page.south east){\includegraphics[scale=0.8, angle=180, origin=c]{img/corner}};
		
		
		% Original pgfornament style:
		%\node[anchor=north west, xshift=0.5cm, yshift=-0.5cm,inner sep=0] at (current page.north west){\pgfornament[scale=0.5]{41}};
		%\node[anchor=north east, xshift=-0.5cm, yshift=-0.5cm, inner sep=0] at (current page.north east){\pgfornament[scale=0.5,symmetry=v]{41}};
		%\node[anchor=south west, xshift=0.5cm, yshift=0.5cm, inner sep=0] at (current page.south west){\pgfornament[scale=0.5,symmetry=h]{41}};
		%\node[anchor=south east, xshift=-0.5cm, yshift=0.5cm, inner sep=0] at (current page.south east){\pgfornament[scale=0.5,symmetry=c]{41}};
		
		% Background as an image, for example a parchment like texture
		%\tikz[remember picture, overlay] \node[opacity=0.5, inner sep=0pt] at (current page.center){\includegraphics[width=\paperwidth,height=\paperheight]{example-image}};
	\end{tikzpicture}
	
	\begin{longtable}{p{1cm}p{10cm}}
		1 -- 2 & Testfigur\\
		3 -- 4 & Test2 blarg\\
		5 -- 8 & Test1\\
		
	\end{longtable}
	
\end{document}